\section{Signed Circulants and Flux Coordinates}

\subsection{Signed adjacency matrices and switching}

For an integer \(n\ge 8\), let \(G_n=C_n(1,2)\) have vertex set \(Z/nZ\) and edges

\begin{lstlisting}[language={}]
{i,i+1}, {i,i+2}    (i in Z/nZ).
\end{lstlisting}

A signing is a map \(\sigma:E(G_n)->{+1,-1}\). Its signed adjacency matrix
\(A_\sigma\) is the real symmetric matrix whose \(ij\) entry is the sign of
\({i,j}\) when this is an edge and zero otherwise. We write

\begin{lstlisting}[language={}]
rho(A_sigma)=max{|lambda|: lambda in spec(A_sigma)}.
\end{lstlisting}

For a vertex-sign function \(d:Z/nZ->{+1,-1}\), let \(D=\operatorname{diag}(d_i)\). Replacing
\(\sigma_ij\) by \(d_i \sigma_ij d_j\) replaces \(A_\sigma\) by \(D A_\sigma D\) and
therefore preserves the spectrum. This operation is switching.

Because \(G_n\) is connected and has \(2n\) edges, its cycle space has dimension
\(2n-n+1=n+1\). Hence there are exactly \(2^{n+1}\) switching classes. We use a
cycle-flux coordinate system adapted to the Hamilton cycle of step-one edges.

\subsection{Hamilton gauge, triangle flux, and quadrilateral flux}

Let \(a_i\) be the sign of the step-one edge \({i,i+1}\) and \(b_i\) the sign of
the step-two edge \({i,i+2}\). Define the triangle flux and Hamilton holonomy by

\begin{lstlisting}[language={}]
tau_i=a_i a_(i+1) b_i,
alpha=product_(i=0)^(n-1) a_i.
\end{lstlisting}

The \(n\) triangle cycles together with the step-one Hamilton cycle form a cycle
basis, so \((\tau_0,...,\tau_{n-1},\alpha)\) determines the switching class.

Switch so that all step-one edges are positive except possibly the edge
crossing the cut from \(n-1\) to \(0\). Equivalently, extend a vector to the
integer lattice with twisted boundary condition

\begin{lstlisting}[language={}]
x_(i+n)=alpha x_i.
\end{lstlisting}

The signed operator then has the local form

\begin{lstlisting}[language={}]
(A_tau x)_i=x_(i-1)+x_(i+1)+tau_(i-2)x_(i-2)+tau_i x_(i+2).       (2.1)
\end{lstlisting}

The distinguished local quadrilateral flux word is

\begin{lstlisting}[language={}]
Q_i=tau_i tau_(i+1).                                             (2.2)
\end{lstlisting}

For a \(p\)-periodic word \(\tau\), equation (2.2) implies
\(\prod_{i=0}^{p-1} Q_i=1\). Conversely, any word
\(Q \in {+1,-1}^p\) with product \(+1\) has exactly two periodic lifts, obtained by
choosing \(\tau_0=+1\) or \(-1\) and recursing via
\(\tau_{i+1}=Q_i \tau_i\).

We call

\begin{lstlisting}[language={}]
D(Q)={i:Q_i=+1}
\end{lstlisting}

the defect set and write \(d=|D(Q)|\). We also use the cyclic local statistics

\begin{lstlisting}[language={}]
a=#{i:Q_i=Q_(i+1)=+1},
b=#{i:Q_i=Q_(i+2)=+1}.                                          (2.3)
\end{lstlisting}

When \(n=8\), the graph contains two additional step-two quadrilaterals not
among the local cycles encoded by (2.2). This causes no loss of finite
information: the cycle-basis coordinates \((\tau,\alpha)\), rather than \(Q\) alone,
parametrize every switching class and are the coordinates used in the
all-signings enumeration.

\subsection{Floquet matrices}

Suppose \(\tau\) has displayed period \(p\). Write an integer index as \(mp+r\),
where \(0\le r<p\), and seek Bloch vectors of the form

\begin{lstlisting}[language={}]
x_(mp+r)=z^m v_r.
\end{lstlisting}

Substitution in (2.1) gives a \(p x p\) Laurent matrix \(H_\tau(z)\). More
explicitly, for each transition from output residue \(r\) to an absolute source
index \(j\), write \(j=mp+s\), \(0\le s<p\), and add the transition coefficient times
\(z^m\) to entry \((r,s)\). This convention automatically adds coefficients when
different transitions collide in short cells.

The identity

\begin{lstlisting}[language={}]
H_tau(z)^T=H_tau(z^(-1))                                        (2.4)
\end{lstlisting}

follows by reversing each undirected transition. Thus \(H_\tau(z)\) is Hermitian
on the unit circle.

For an infinite periodic phase, define

\begin{lstlisting}[language={}]
R(Q)=sup_(|z|=1) rho(H_tau(z))^2.                                (2.5)
\end{lstlisting}

The two lifts of \(Q\) have the same value, so this notation is unambiguous.
Notice that \(R(Q)\) is always a \textbf{squared} spectral radius.

For every explicit fiber certificate below we use the canonical lift
\(\tau_0=+1\), \(\tau_{i+1}=Q_i \tau_i\), and the ordered fiber basis
\((v_0,...,v_{p-1})\). The other lift is obtained by the diagonal conjugacy and
fiber reparametrization in Lemma 2.1, so it has the same squared spectral
conclusions.

Whenever an explicit matrix representative is required, we write
\(H_Q(z):=H_{\tau^can}(z)\) for the fiber formed from this canonical lift and
basis. The notation \(R(Q)\) and the moments \(M_k(Q)\) remain lift-invariant.

If a finite graph has order \(n=pL\) and Hamilton holonomy \(\alpha\), its cell
sequence satisfies \(u_{m+L}=\alpha u_m\). The unitary cell shift has eigenvalues
exactly

\begin{lstlisting}[language={}]
z^L=alpha,                                                       (2.6)
\end{lstlisting}

and the finite matrix decomposes as

\begin{lstlisting}[language={}]
A_(pL,alpha)  ~=  direct_sum_(z^L=alpha) H_tau(z).               (2.7)
\end{lstlisting}

Equation (2.6) is a discrete finite set. Equation \(|z|=1\) in (2.5) describes
the infinite-volume spectrum. We will not substitute one for the other.

\subsection{Operator equivalences}

We record the equivalences used in classifying periodic phases.

\textbf{Lemma 2.1 (translation, reflection, and edge-sign negation).} Translating \(\tau\),
reflecting it by \(\tau_i -> \tau_{-i-2}\), or replacing \(\tau\) by \(-\tau\) preserves
\(R(Q)\). On flux words reflection is \(Q_i -> Q_{-i-3}\).

\textbf{Proof.} Let \((T_r x)_i=x_{i-r}\) and \((Jx)_i=x_{-i}\). Direct substitution in
(2.1) gives

\begin{lstlisting}[language={}]
T_r A_tau T_r^(-1)=A_(translated tau),
J A_tau J=A_(reflected tau).
\end{lstlisting}

Both are unitary conjugacies. For global edge-sign negation let
\((Dx)_i=(-1)^i x_i\). The
endpoints of a step-one edge have opposite \(D\) signs while those of a step-two
edge have equal signs. Therefore

\begin{lstlisting}[language={}]
A_(-tau)=-D A_tau D.
\end{lstlisting}

The spectrum is negated and its square is unchanged. In an odd displayed cell
the fiber coordinate changes from \(z\) to \(-z\), but the full unit circle is
preserved. \(\square\)

\textbf{Lemma 2.2 (zone folding).}  If \(\tau\) has primitive period \(q\) and is
displayed in a repeated cell \(p=mq\), then

\begin{lstlisting}[language={}]
H_p(z) ~= direct_sum_(w^m=z) H_q(w).                            (2.8)
\end{lstlisting}

In particular repetition of the unit cell preserves \(R(Q)\).

\textbf{Proof.} Write a residue in the repeated cell as \(r+kq\), where
\(0\le r<q\) and \(0\le k<m\). Decompose the \(mq\)-dimensional fiber according to the
unitary internal translation by \(q\). On its \(w\)-eigenspace vectors have the
form

\begin{lstlisting}[language={}]
x_(r+kq)=w^k v_r.
\end{lstlisting}

The repeated boundary condition is equivalent to \(w^m=z\). Since all
coefficients in (2.1) are \(q\)-periodic, every transition factors out \(w^k\)
and the remaining action on \(v\) is exactly \(H_q(w)\). The \(m\) roots of
\(w^m=z\) give mutually orthogonal eigenspaces and their dimensions sum to
\(mq\). This proves (2.8), including multiplicities. \(\square\)

\subsection{Closed-walk moments}

For a periodic phase define

\begin{lstlisting}[language={}]
M_k(Q)=CT_z tr(H_Q(z)^(2k)),
F_k(Q)=M_(k+1)(Q)-8M_k(Q).                                      (2.9)
\end{lstlisting}

Constant-term extraction equals normalized integration around the unit circle:

\begin{lstlisting}[language={}]
M_k(Q)=(1/(2pi)) integral_0^(2pi)
       sum_j lambda_j(e^(i theta))^(2k) dtheta.                  (2.10)
\end{lstlisting}

It is also the signed sum of closed walks of length \(2k\) per period cell.

\textbf{Lemma 2.3 (moment barrier).} If \(R(Q)\le 8\), then
\(M_{k+1}(Q)\le 8M_k(Q)\) for every \(k\ge 1\). Equivalently,

\begin{lstlisting}[language={}]
F_k(Q)>0  ==>  R(Q)>8.                                          (2.11)
\end{lstlisting}

\textbf{Proof.} Under \(R(Q)\le 8\), every
\(y_j(\theta)=\lambda_j(e^{i \theta})^2\) lies in \([0,8]\). Hence
\(y_j^{k+1}\le 8y_j^k\). Sum over \(j\) and integrate using (2.10). \(\square\)

We emphasize that (2.11) is one-way. A nonpositive excess, or any finite list
of nonpositive excesses, does not prove \(R(Q)\le 8\).

\subsection{Distinguished constants}

For even \(n\ge 8\), the conjectured threshold inherited from \cite{Suvagiya2026Signed} is

\begin{lstlisting}[language={}]
rho_-(n)=2 sqrt(cos^2(pi/n)+cos^2(2pi/n)),
rho_-(n)^2=4+2cos(2pi/n)+2cos(4pi/n).                            (2.12)
\end{lstlisting}

The target period-eight squared spectral edge is

\begin{lstlisting}[language={}]
eta=4+sqrt(10+2sqrt(5)),
rho_*=sqrt(eta).                                                  (2.13)
\end{lstlisting}

These quantities play different roles: \(\rho_-(n)\) is the finite conjectured
lower bound, while \(\eta\) is the exact infinite-volume squared radius of the
counterexample phase.
